\documentclass[12pt,letterpaper]{article}
\usepackage{graphicx,textcomp}
\usepackage{natbib}
\usepackage{setspace}
\usepackage{fullpage}
\usepackage{color}
\usepackage[reqno]{amsmath}
\usepackage{amsthm}
\usepackage{fancyvrb}
\usepackage{amssymb,enumerate}
\usepackage[all]{xy}
\usepackage{endnotes}
\usepackage{lscape}
\newtheorem{com}{Comment}
\usepackage{float}
\usepackage{hyperref}
\newtheorem{lem} {Lemma}
\newtheorem{prop}{Proposition}
\newtheorem{thm}{Theorem}
\newtheorem{defn}{Definition}
\newtheorem{cor}{Corollary}
\newtheorem{obs}{Observation}
\usepackage[compact]{titlesec}
\usepackage{dcolumn}
\usepackage{tikz}
\usetikzlibrary{arrows}
\usepackage{multirow}
\usepackage{xcolor}
\newcolumntype{.}{D{.}{.}{-1}}
\newcolumntype{d}[1]{D{.}{.}{#1}}
\definecolor{light-gray}{gray}{0.65}
\usepackage{url}
\usepackage{listings}
\usepackage{color}

\definecolor{codegreen}{rgb}{0,0.6,0}
\definecolor{codegray}{rgb}{0.5,0.5,0.5}
\definecolor{codepurple}{rgb}{0.58,0,0.82}
\definecolor{backcolour}{rgb}{0.95,0.95,0.92}

\lstdefinestyle{mystyle}{
	backgroundcolor=\color{backcolour},   
	commentstyle=\color{codegreen},
	keywordstyle=\color{magenta},
	numberstyle=\tiny\color{codegray},
	stringstyle=\color{codepurple},
	basicstyle=\footnotesize,
	breakatwhitespace=false,         
	breaklines=true,                 
	captionpos=b,                    
	keepspaces=true,                 
	numbers=left,                    
	numbersep=5pt,                  
	showspaces=false,                
	showstringspaces=false,
	showtabs=false,                  
	tabsize=2
}
\lstset{style=mystyle}
\newcommand{\Sref}[1]{Section~\ref{#1}}
\newtheorem{hyp}{Hypothesis}


\title{Problem Set 4}
\date{Due: November 27, 2025}
\author{Applied Stats/Quant Methods 1}


\begin{document}
	\maketitle
	\section*{Instructions}
	\begin{itemize}
		\item Please show your work! You may lose points by simply writing in the answer. If the problem requires you to execute commands in \texttt{R}, please include the code you used to get your answers. Please also include the \texttt{.R} file that contains your code. If you are not sure if work needs to be shown for a particular problem, please ask.
		\item Your homework should be submitted electronically on GitHub.
		\item This problem set is due before 23:59 on Thursday November 27, 2025. No late assignments will be accepted.
	\end{itemize}



	\vspace{.5cm}
\section*{Question 1: Economics}
\vspace{.25cm}
\noindent 	
In this question, use the \texttt{prestige} dataset in the \texttt{car} library. First, run the following commands:

\begin{verbatim}
install.packages(car)
library(car)
data(Prestige)
help(Prestige)
\end{verbatim} 


\noindent We would like to study whether individuals with higher levels of income have more prestigious jobs. Moreover, we would like to study whether professionals have more prestigious jobs than blue and white collar workers.

\newpage
\begin{enumerate}
	
	\item [(a)]
	Create a new variable \texttt{professional} by recoding the variable \texttt{type} so that professionals are coded as $1$, and blue and white collar workers are coded as $0$ (Hint: \texttt{ifelse}).

	\vspace{.5cm}

	\lstinputlisting[language=R, firstline=40, lastline=41]{PS04_answersKC.R}
	
	\item [(b)]
	Run a linear model with \texttt{prestige} as an outcome and \texttt{income}, \texttt{professional}, and the interaction of the two as predictors (Note: this is a continuous $\times$ dummy interaction.)

	\vspace{.5cm}

	\lstinputlisting[language=R, firstline=43, lastline=44]{PS04_answersKC.R}
	\begin{verbatim}
Call:
lm(formula = prestige ~ income + professional + income:professional, 
	data = Prestige)
		
Coefficients:
		(Intercept)               income         professional  
		  21.142259             0.003171            37.781280  
income:professional  
	      -0.002326
	\end{verbatim}
	
	\item [(c)]
	Write the prediction equation based on the result.

	\vspace{.5cm}

	\noindent Prestige = 21.142259 + 0.003171 * income + 37.781280  * professional + -0.002326 * (income * professional)
	
	\item [(d)]
	Interpret the coefficient for \texttt{income}.

	\vspace{.5cm}

	\noindent For every one unit increase in income there is a 0.003171 increase in prestige on average. However, this interpretation is only true if the dummy variable professional equals 0. Otherwise, the coefficient for income should be interpreted with the interaction effect.

	In this case, we have a coefficient of -0.002326 for interaction, which when added to the income coefficient 0.003171 results in 0.000845. Therefore, if the dummy variable equals 1, we should say that on average for every one unit increase in income there is a 0.000845 increase in prestige.
	
	\newpage
	\item [(e)]
	Interpret the coefficient for \texttt{professional}.
	
	\vspace{.1cm}

	\noindent For a one unit increase in professional (which represents the change from a blue or white collar job to a professional job) there is a 37.7812800 increase in prestige on average if income = 0. Otherwise, the coefficient for professional should be interpreted with the interaction effect and the effect will vary depending on the value of income.
	
	\item [(f)]
	What is the effect of a \$1,000 increase in income on prestige score for professional occupations? In other words, we are interested in the marginal effect of income when the variable \texttt{professional} takes the value of $1$. Calculate the change in $\hat{y}$ associated with a \$1,000 increase in income based on your answer for (c).
	
	\vspace{.1cm}
	
	\noindent Predicted Prestige = 21.1422589 + 0.0031709 * income + 37.7812800  * professional + (-0.0023257) * (income * professional)
	
	\noindent In this example, the marginal effect equals the coefficient for income + the interaction coefficient * 1,000 (because we are looking at the impact of a 1,000 dollar increase in income). By including the interaction coefficient (multiplied by 1), we are holding the dummy variable at 1 because a dummy variable of 0 would have resulted in an interaction effect of 0. So to find the change in y for a 1,000 dollar increase in income, we do (0.0031709 - 0.0023257) * 1,000, which equals 0.8452. So, for a 1,000 dollar increase in income for professionals, we would expect a 0.8452 increase in prestige.
	
	\lstinputlisting[language=R, firstline=47, lastline=47]{PS04_answersKC.R}
	\begin{verbatim}
	[1] 0.8452
	\end{verbatim}
	
	\item [(g)]
	What is the effect of changing one's occupations from non-professional to professional when her income is \$6,000? We are interested in the marginal effect of professional jobs when the variable \texttt{income} takes the value of $6,000$. Calculate the change in $\hat{y}$ based on your answer for (c).
	
	\vspace{.1cm}
	
	\noindent Predicted Prestige = 21.1422589 + 0.0031709 * income + 37.7812800  * professional + (-0.0023257) * (income * professional)
	
	\noindent To find the marginal effect of professional jobs when the variable income takes the value of 6,000, I added the coefficient for professional jobs to the interaction effect * 6000. When substituting in the numbers the equation was 37.7812800 + (-0.0023257 * 6000), which equals 23.82708. So, if income takes the value of 6,000, we would expect the average difference between prestige for a blue/white collar job and prestige for a professional job to be 23.82708.
	
	\lstinputlisting[language=R, firstline=49, lastline=49]{PS04_answersKC.R}
	
	\begin{verbatim}
	[1] 23.82708
	\end{verbatim} 
	

	


	
	
	
\end{enumerate}

\vspace{1cm}

\section*{Question 2: Political Science}
\vspace{.25cm}
\noindent 	Researchers are interested in learning the effect of all of those yard signs on voting preferences.\footnote{Donald P. Green, Jonathan	S. Krasno, Alexander Coppock, Benjamin D. Farrer,	Brandon Lenoir, Joshua N. Zingher. 2016. ``The effects of lawn signs on vote outcomes: Results from four randomized field experiments.'' Electoral Studies 41: 143-150. } Working with a campaign in Fairfax County, Virginia, 131 precincts were randomly divided into a treatment and control group. In 30 precincts, signs were posted around the precinct that read, ``For Sale: Terry McAuliffe. Don't Sellout Virgina on November 5.'' \\

Below is the result of a regression with two variables and a constant.  The dependent variable is the proportion of the vote that went to McAuliff's opponent Ken Cuccinelli. The first variable indicates whether a precinct was randomly assigned to have the sign against McAuliffe posted. The second variable indicates
a precinct that was adjacent to a precinct in the treatment group (since people in those precincts might be exposed to the signs).  \\

\vspace{.5cm}
\begin{table}[!htbp]
	\centering 
	\textbf{Impact of lawn signs on vote share}\\
	\begin{tabular}{@{\extracolsep{5pt}}lccc} 
		\\[-1.8ex] 
		\hline \\[-1.8ex]
		Precinct assigned lawn signs  (n=30)  & 0.042\\
		& (0.016) \\
		Precinct adjacent to lawn signs (n=76) & 0.042 \\
		&  (0.013) \\
		Constant  & 0.302\\
		& (0.011)
		\\
		\hline \\
	\end{tabular}\\
	\footnotesize{\textit{Notes:} $R^2$=0.094, N=131}
\end{table}

\vspace{.5cm}
\begin{enumerate}
	\item [(a)] Use the results from a linear regression to determine whether having these yard signs in a precinct affects vote share (e.g., conduct a hypothesis test with $\alpha = .05$).
	\vspace{.5cm}
	
	\noindent \textbf{State Assumptions}
	
	\begin{itemize}
	\item n is greater than or equal to 30
	\item The data is continuous
	\item The sample was random
	\end{itemize}
	
	\newpage
	
	\noindent \textbf{State Hypotheses}
	\begin{itemize}
	\item Null Hypothesis: $$\beta1 = 0$$
	\item Alternative Hypothesis: $$\beta1 \neq 0$$
	\end{itemize}
	
	\noindent \textbf{Calculate the Test Statistic}
	
	\noindent (Observed slope - Slope under the null hypothesis)/Standard error
	
	\noindent The sample slope is 0.042. The slope under the null hypothesis is 0. The sample standard error is 0.016. Therefore, the test statistic = (0.042 - 0)/0.016 or 0.042/0.016, which equals 2.625.
	
	\lstinputlisting[language=R, firstline=54, lastline=55]{PS04_answersKC.R}
	
	\begin{verbatim}
	[1] 2.625
	\end{verbatim}
	
	\noindent \textbf{Calculate the p-value}
	
	\noindent To calculate the p-value, I needed the degrees of freedom, which I calculated using the formula n - p - 1, where n is the sample size and p is the number of predictor variables in the model. In this example, n = 30 (I used the n for each group not the combined N) and p = 2 (because there are 2 predictor variables). 30 - 2 - 1 = 27
	
	\lstinputlisting[language=R, firstline=57, lastline=60]{PS04_answersKC.R}
	
	\begin{verbatim}
	[1] 27
	\end{verbatim}
	
	\noindent After finding the degrees of freedom, I calculated the p-value by using the pt function and inputting the values I found for the test statistic and the degrees of freedom. I set lower.tail equal to FALSE and multiplied the result by 2 because this is a two-tailed t-test. This is because we do not care if our slope is greater than or less than 0, we only care whether or not it is significantly different from 0.
	
	\lstinputlisting[language=R, firstline=61, lastline=62]{PS04_answersKC.R}
	
	\begin{verbatim}
	[1] 0.01409096
	\end{verbatim}
	
	\noindent I got a p-value of 0.01409096, which is less than the critical value of 0.05. Therefore, we can reject the null hypothesis that Beta 1 (the slope for precincts assigned lawn signs) = 0
	
	\item [(b)]  Use the results to determine whether being
	next to precincts with these yard signs affects vote
	share (e.g., conduct a hypothesis test with $\alpha = .05$).
	
	\noindent \textbf{State Assumptions}
	
	\begin{itemize}
		\item n is greater than or equal to 30
		\item The data is continuous
		\item The sample was random
	\end{itemize}
	
	\noindent \textbf{State Hypotheses}
	\begin{itemize}
		\item Null Hypothesis: $$\beta2 = 0$$
		\item Alternative Hypothesis: $$\beta2 \neq 0$$
	\end{itemize}
	
	\noindent \textbf{Calculate the Test Statistic}
	
	\noindent (Observed slope - Slope under the null hypothesis)/Standard error
	
	\noindent The sample slope is 0.042. The slope under the null hypothesis is 0. The standard error is 0.013. Therefore, the test statistic = (0.042 - 0)/0.013 or 0.042/0.013, which equals 3.230769.
	
	\lstinputlisting[language=R, firstline=64, lastline=65]{PS04_answersKC.R}
	
	\begin{verbatim}
	[1] 3.230769
	\end{verbatim}
	
	\noindent \textbf{Calculate the p-value}
	
	\noindent To calculate the p-value, I needed the degrees of freedom, which I calculated using the formula n - p - 1, where n is the sample size and p is the number of predictor variables in the model. In this example, n = 76 (I used the n for each group not the combined N) and p = 2 (because there are 2 predictor variables). 76 - 2 - 1 = 73
	
	\lstinputlisting[language=R, firstline=67, lastline=69]{PS04_answersKC.R}
	
	\begin{verbatim}
	[1] 73
	\end{verbatim}
	
	After finding the degrees of freedom, I calculated the p-value by using the pt function and inputting the values I found for the test statistic and the degrees of freedom. I set lower.tail equal to FALSE and multiplied the result by 2 because this is a two-tailed t-test. This is because we do not care if our slope is greater than or less than 0, we only care whether or not it is significantly different from 0.
	
	\lstinputlisting[language=R, firstline=71, lastline=72]{PS04_answersKC.R}
	
	\begin{verbatim}
	[1] 0.001852655
	\end{verbatim}
	
	\noindent I got a p-value of 0.001852655, which is less than the critical value of 0.05. Therefore, we can reject the null hypothesis that Beta 2 (the slope for precincts adjacent to lawn signs) = 0
	
	\item [(c)] Interpret the coefficient for the constant term substantively.
	\vspace{0.5cm}
	
	\noindent The constant term suggests that if there were no sign against McAuliff posted, a proportion of 0.302 or 30.2 percent of the vote share would go to McAuliff's opponent Ken Cuccinelli.
	\vspace{.5cm}
	
	\item [(d)] Evaluate the model fit for this regression.  What does this	tell us about the importance of yard signs versus other factors that are not modeled?
	
	\vspace{.5cm}
	
	\noindent The R squared value of 0.094 suggests that the explanatory variable of lawn signs explains very little of the variation in the outcome of vote share. The low R squared value indicates an extremely poor model fit. This suggests that other factors that are not modeled may be more effective in explaining variation in the outcome.
	
\end{enumerate}  


\end{document}
